\documentclass{article}
\usepackage{amsmath,amssymb,amsthm}
\usepackage{algorithm}
\usepackage[noend]{algpseudocode}
\usepackage{enumitem}
\usepackage{hyperref}
\newtheorem{theorem}{Theorem}
\newtheorem{lemma}{Lemma}
\newtheorem{assumption}{Assumption}
\newtheorem{corollary}{Corollary}
\newcommand{\E}{\mathbb{E}}
\newcommand{\R}{\mathbb{R}}

\begin{document}

\section*{Top-$K$ Stochastic Gradient Descent (Top-$K$ SGD)}

\section*{Mathematical Formulation}
Let $f:\R^{d}\rightarrow\R$ be a (possibly non‑convex) differentiable loss.  
At iteration $t$ we have the current iterate $w_{t}\in\R^{d}$ and a stochastic gradient
\[
g_{t}:=\nabla f(w_{t};\xi_{t})\in\R^{d},
\]
where $\xi_{t}$ denotes the random data sample (or mini‑batch).  

Define the \emph{Top-$K$ sparsification operator}
\[
\mathcal{S}_{K}(v):=\left\{
\begin{array}{ll}
v_{i}, & i\in\mathcal{I}_{K}(v),\\[2mm]
0, & \text{otherwise},
\end{array}\right.
\]
where $\mathcal{I}_{K}(v)$ is the index set of the $K$ components of $v$ having the largest absolute values (ties are broken uniformly at random).  
The sparse gradient transmitted to the parameter server is
\[
\tilde g_{t}:=\mathcal{S}_{K}(g_{t}) .
\]

The update rule of Top‑$K$ SGD is
\[
\boxed{\,w_{t+1}=w_{t}-\eta_{t}\,\tilde g_{t}\,} \tag{1}
\]
with learning‑rate schedule $\{\eta_{t}\}_{t\ge 0}>0$ and sparsity level $K\in\{1,\dots ,d\}$ (or equivalently a compression ratio $\tau=K/d$).

\section*{Pseudocode}
\begin{algorithm}[H]
\caption{Top-$K$ SGD}
\begin{algorithmic}[1]
\State \textbf{Input:} initial point $w_{0}\in\R^{d}$, step sizes $\{\eta_{t}\}$, sparsity $K$.
\For{$t=0,1,2,\dots,T-1$}
    \State Sample mini‑batch $\xi_{t}$ and compute stochastic gradient $g_{t}\gets\nabla f(w_{t};\xi_{t})$.
    \State Determine index set $\mathcal{I}_{K}(g_{t})$ of the $K$ largest $|g_{t,i}|$.
    \State Form sparse gradient $\tilde g_{t}\gets\mathcal{S}_{K}(g_{t})$.
    \State Update parameters: $w_{t+1}\gets w_{t}-\eta_{t}\tilde g_{t}$.
\EndFor
\State \textbf{Output:} $w_{T}$ (or the average $\bar w_{T}:=\tfrac1T\sum_{t=0}^{T-1}w_{t}$).
\end{algorithmic}
\end{algorithm}

\section*{Dry Run Example}
Consider the scalar quadratic $f(w)=(w-3)^{2}$, thus $\nabla f(w)=2(w-3)$.  
Take $w_{0}=0$, constant step size $\eta=0.1$, and $K=1$ (the only coordinate is kept).

\begin{center}
\begin{tabular}{c|c|c|c}
$t$ & $g_{t}=2(w_{t}-3)$ & $\tilde g_{t}=g_{t}$ (Top‑$1$ keeps it) & $w_{t+1}=w_{t}-\eta\tilde g_{t}$ \\ \hline
0 & $2(0-3)=-6$ & $-6$ & $0-0.1(-6)=0.6$ \\
1 & $2(0.6-3)=-4.8$ & $-4.8$ & $0.6-0.1(-4.8)=1.08$ \\
2 & $2(1.08-3)=-3.84$ & $-3.84$ & $1.08-0.1(-3.84)=1.464$ \\
\end{tabular}
\end{center}
The corresponding objective values are $f(0)=9$, $f(0.6)=5.76$, $f(1.08)=3.6864$, $f(1.464)=2.3329$, showing descent.

\vspace{1em}
\section*{Assumptions}
\begin{assumption}[Smoothness]\label{as:smooth}
The loss $f$ has $L$‑Lipschitz continuous gradient:
\[
\|\nabla f(x)-\nabla f(y)\|\le L\|x-y\|,\qquad\forall x,y\in\R^{d}.
\]
\end{assumption}
\begin{assumption}[Unbiased Stochastic Gradient]\label{as:unbiased}
For every $t$, conditioned on $w_{t}$,
\[
\E_{\xi_{t}}[g_{t}\mid w_{t}]=\nabla f(w_{t}),\qquad
\E_{\xi_{t}}\big[\|g_{t}-\nabla f(w_{t})\|^{2}\mid w_{t}\big]\le\sigma^{2}.
\]
\end{assumption}
\begin{assumption}[Top‑$K$ Approximation Error]\label{as:topk}
For any vector $v\in\R^{d}$,
\[
\|v-\mathcal{S}_{K}(v)\|\le \Bigl(1-\frac{K}{d}\Bigr)\|v\|.
\]
Consequently,
\[
\|\nabla f(w_{t})-\tilde g_{t}\|\le\Bigl(1-\frac{K}{d}\Bigr)\|\nabla f(w_{t})\|+ \|g_{t}-\nabla f(w_{t})\|.
\]
\end{assumption}
\begin{assumption}[Step‑size]\label{as:stepsize}
The learning rate is constant $\eta_{t}\equiv\eta\le\frac{1}{L}$.
\end{assumption}

\section*{Main Convergence Theorem}
\begin{theorem}[Convergence of Top‑$K$ SGD]\label{thm:main}
Assume \ref{as:smooth}–\ref{as:stepsize}. Let $\{w_{t}\}_{t=0}^{T}$ be generated by \eqref{1}. Then
\[
\frac{1}{T}\sum_{t=0}^{T-1}\E\big[\|\nabla f(w_{t})\|^{2}\big]
\;\le\;
\underbrace{\frac{2\big(f(w_{0})-f^{*}\big)}{\eta T}}_{\text{optimization term}}
\;+\;
\underbrace{2L\eta\sigma^{2}}_{\text{stochastic noise}}
\;+\;
\underbrace{4L\Bigl(1-\frac{K}{d}\Bigr)^{2}\big(f(w_{0})-f^{*}\big)}_{\text{compression bias}} .
\tag{2}
\]
In particular, choosing $\eta=\Theta(1/\sqrt{T})$ yields
\[
\frac{1}{T}\sum_{t=0}^{T-1}\E\big[\|\nabla f(w_{t})\|^{2}\big]=\mathcal{O}\!\Bigl(\frac{1}{\sqrt{T}}+\Bigl(1-\frac{K}{d}\Bigr)^{2}\Bigr).
\]
Thus Top‑$K$ SGD attains the same $\mathcal{O}(1/\sqrt{T})$ rate as vanilla SGD up to an additive term that vanishes when $K=d$ (i.e., no compression).
\end{theorem}

\section*{Theoretical Properties}
\begin{itemize}[leftmargin=*]
\item \textbf{Stability.} The bound \eqref{2} holds for any $K\ge1$; the compression term scales as $(1-K/d)^{2}$, guaranteeing stability even for aggressive sparsification when $K$ is a constant fraction of $d$.
\item \textbf{Hyper‑parameter Sensitivity.}
    \begin{itemize}
    \item Step size $\eta$ must satisfy $\eta\le 1/L$ (Assumption~\ref{as:stepsize}) for the smoothness inequality to contract.
    \item Larger $K$ reduces the bias term; when $K=d$ the bound collapses to the classical SGD rate.
    \end{itemize}
\item \textbf{Comparison to Baselines.}
    \begin{itemize}
    \item Compared with full‑gradient SGD the extra term $4L(1-K/d)^{2}(f(w_{0})-f^{*})$ quantifies exactly the loss due to sparsification.
    \item When $K\ll d$, the bound remains meaningful because the compression term is independent of $T$ and thus does not affect the asymptotic $\mathcal{O}(1/\sqrt{T})$ decay.
    \end{itemize}
\end{itemize}

\section*{Discussion}
The theorem shows that Top‑$K$ SGD inherits the optimal $\mathcal{O}(1/\sqrt{T})$ convergence rate for non‑convex smooth objectives up to a constant compression penalty. The penalty disappears as the sparsity level $K$ approaches the full dimension, confirming that the algorithm interpolates continuously between extreme compression and the classical SGD regime. The proof technique mirrors the standard smooth‑function descent argument, with an additional decomposition to control the error introduced by the Top‑$K$ operator.

\newpage

\section*{Preliminaries (Definitions and Lemmas)}
\begin{lemma}[Smoothness Inequality]\label{lem:smooth}
If $f$ satisfies Assumption~\ref{as:smooth}, then for any $x,y\in\R^{d}$,
\[
f(y)\le f(x)+\langle\nabla f(x),y-x\rangle+\frac{L}{2}\|y-x\|^{2}.
\]
\end{lemma}
\begin{proof}
Standard consequence of $L$‑Lipschitz gradients; see e.g. Nesterov (2004). \qedhere
\end{proof}

\begin{lemma}[Top‑$K$ Error Bound]\label{lem:topk}
For any $v\in\R^{d}$, the sparsifier $\mathcal{S}_{K}$ satisfies
\[
\|v-\mathcal{S}_{K}(v)\|\le\Bigl(1-\frac{K}{d}\Bigr)\|v\|.
\]
\end{lemma}
\begin{proof}
Let $\mathcal{I}_{K}(v)$ be the index set of the $K$ largest absolute components.
By definition the $d-K$ discarded components have magnitude at most the $K$‑th largest, hence
\[
\sum_{i\notin\mathcal{I}_{K}}v_{i}^{2}\le\frac{d-K}{K}\sum_{i\in\mathcal{I}_{K}}v_{i}^{2}
\]
which yields the claimed inequality after elementary algebra. \qedhere
\end{proof}

\section*{Main Proof (Step‑by‑Step Derivation)}
We prove Theorem~\ref{thm:main}. Throughout the proof, expectations are taken w.r.t. all randomness up to time $t$ (including the stochastic sample $\xi_{t}$ and the random tie‑breaking in $\mathcal{S}_{K}$).

\noindent\textbf{Step 1 (Smoothness descent).}  
Applying Lemma~\ref{lem:smooth} with $x=w_{t}$ and $y=w_{t+1}=w_{t}-\eta\tilde g_{t}$ gives
\begin{align}
f(w_{t+1})
&\le f(w_{t})-\eta\langle\nabla f(w_{t}),\tilde g_{t}\rangle
   +\frac{L\eta^{2}}{2}\|\tilde g_{t}\|^{2}. \tag{3}
\end{align}
[Step 1]

\noindent\textbf{Step 2 (Decompose the inner product).}  
Write
\[
\langle\nabla f(w_{t}),\tilde g_{t}\rangle
= \|\nabla f(w_{t})\|^{2}
   -\langle\nabla f(w_{t}),\nabla f(w_{t})-\tilde g_{t}\rangle.
\]
Using Cauchy–Schwarz,
\[
\langle\nabla f(w_{t}),\tilde g_{t}\rangle
\ge \|\nabla f(w_{t})\|^{2}
   -\|\nabla f(w_{t})\|\;\|\nabla f(w_{t})-\tilde g_{t}\|. \tag{4}
\]
[Step 2]

\noindent\textbf{Step 3 (Bound the compression error).}  
From Assumption~\ref{as:topk} and the unbiasedness of $g_{t}$,
\[
\|\nabla f(w_{t})-\tilde g_{t}\|
\le\Bigl(1-\frac{K}{d}\Bigr)\|\nabla f(w_{t})\|
    +\|g_{t}-\nabla f(w_{t})\|. \tag{5}
\]
[Step 3]

\noindent\textbf{Step 4 (Plug (4)–(5) into (3)).}  
Using (4) and (5) in (3) yields
\begin{align}
f(w_{t+1})
&\le f(w_{t})-\eta\|\nabla f(w_{t})\|^{2}
   +\eta\|\nabla f(w_{t})\|
      \Bigl[\Bigl(1-\frac{K}{d}\Bigr)\|\nabla f(w_{t})\|
            +\|g_{t}-\nabla f(w_{t})\|\Bigr] \nonumber\\
&\quad+\frac{L\eta^{2}}{2}\|\tilde g_{t}\|^{2}. \tag{6}
\end{align}
[Step 4]

\noindent\textbf{Step 5 (Upper bound $\|\tilde g_{t}\|^{2}$).}  
Since $\tilde g_{t}$ is a sub‑vector of $g_{t}$,
\[
\|\tilde g_{t}\|^{2}\le\|g_{t}\|^{2}
   =\|\nabla f(w_{t})\|^{2}
    +\|g_{t}-\nabla f(w_{t})\|^{2}
    +2\langle\nabla f(w_{t}),g_{t}-\nabla f(w_{t})\rangle.
\]
Taking conditional expectation and using $\E[g_{t}\mid w_{t}]=\nabla f(w_{t})$ eliminates the cross term, so
\[
\E\big[\|\tilde g_{t}\|^{2}\mid w_{t}\big]
\le \|\nabla f(w_{t})\|^{2}+ \E\big[\|g_{t}-\nabla f(w_{t})\|^{2}\mid w_{t}\big]
\le \|\nabla f(w_{t})\|^{2}+ \sigma^{2}. \tag{7}
\]
[Step 5]

\noindent\textbf{Step 6 (Take expectation of (6)).}  
Condition on $w_{t}$ and apply (7):
\begin{align}
\E\!\big[f(w_{t+1})\mid w_{t}\big]
&\le f(w_{t})-\eta\|\nabla f(w_{t})\|^{2}
   +\eta\Bigl(1-\frac{K}{d}\Bigr)\|\nabla f(w_{t})\|^{2}
   +\eta\|\nabla f(w_{t})\|\E\!\big[\|g_{t}-\nabla f(w_{t})\|\mid w_{t}\big] \nonumber\\
&\quad+\frac{L\eta^{2}}{2}\big(\|\nabla f(w_{t})\|^{2}+\sigma^{2}\big). \tag{8}
\end{align}
[Step 6]

\noindent\textbf{Step 7 (Bound the mixed term $\|\nabla f\|\E\|g-\nabla f\|$).}  
By Jensen’s inequality,
\[
\E\!\big[\|g_{t}-\nabla f(w_{t})\|\mid w_{t}\big]
\le\big(\E\!\big[\|g_{t}-\nabla f(w_{t})\|^{2}\mid w_{t}\big]\big)^{1/2}
\le\sigma.
\]
Thus
\[
\eta\|\nabla f(w_{t})\|\E\!\big[\|g_{t}-\nabla f(w_{t})\|\mid w_{t}\big]
\le \eta\sigma\|\nabla f(w_{t})\|
\le \frac{\eta}{2}\|\nabla f(w_{t})\|^{2}+\frac{\eta\sigma^{2}}{2},
\]
where the last inequality follows from $ab\le \frac{a^{2}}{2}+\frac{b^{2}}{2}$ with $a=\|\nabla f(w_{t})\|$ and $b=\sigma$. \tag{9}
[Step 7]

\noindent\textbf{Step 8 (Collect terms).}  
Insert (9) into (8) and simplify, using $\eta\le 1/L$ (Assumption~\ref{as:stepsize}) to keep the coefficient of $\|\nabla f(w_{t})\|^{2}$ non‑negative:
\begin{align}
\E\!\big[f(w_{t+1})\mid w_{t}\big]
&\le f(w_{t})
   -\underbrace{\Bigl(\eta-\eta\Bigl(1-\frac{K}{d}\Bigr)-\frac{\eta}{2}-\frac{L\eta^{2}}{2}\Bigr)}_{\displaystyle\alpha}
     \|\nabla f(w_{t})\|^{2}
   +\underbrace{\Bigl(\frac{\eta\sigma^{2}}{2}+\frac{L\eta^{2}\sigma^{2}}{2}\Bigr)}_{\displaystyle\beta}. \tag{10}
\end{align}
Since $\eta\le 1/L$, we have $L\eta^{2}\le\eta$, and therefore
\[
\alpha \;\ge\; \frac{\eta}{2}\Bigl(\frac{K}{d}\Bigr). 
\]
Consequently,
\[
\E\!\big[f(w_{t+1})\mid w_{t}\big]
\le f(w_{t})-\frac{\eta K}{2d}\|\nabla f(w_{t})\|^{2}+\beta. \tag{11}
\]
[Step 8]

\noindent\textbf{Step 9 (Unroll the recursion).}  
Take total expectation and sum (11) over $t=0,\dots,T-1$:
\[
\E[f(w_{T})]
\le f(w_{0})
 -\frac{\eta K}{2d}\sum_{t=0}^{T-1}\E\big[\|\nabla f(w_{t})\|^{2}\big]
 +T\beta .
\]
Rearranging,
\[
\frac{1}{T}\sum_{t=0}^{T-1}\E\big[\|\nabla f(w_{t})\|^{2}\big]
\le \frac{2d}{\eta K}\,\frac{f(w_{0})-\E[f(w_{T})]}{T}
     +\frac{2d}{K}\,\beta . \tag{12}
\]
Since $f(w_{T})\ge f^{*}$ we can replace $\E[f(w_{T})]$ by $f^{*}$, and recall $\beta=\frac{\eta\sigma^{2}}{2}\bigl(1+L\eta\bigr)\le\eta\sigma^{2}$ (using $\eta\le1/L$). Plugging these into (12) yields
\[
\frac{1}{T}\sum_{t=0}^{T-1}\E\big[\|\nabla f(w_{t})\|^{2}\big]
\le \frac{2d}{\eta K}\frac{f(w_{0})-f^{*}}{T}
   +\frac{2d}{K}\,\eta\sigma^{2}. \tag{13}
\]
[Step 9]

\noindent\textbf{Step 10 (Express the bound in the original form).}  
Observe that $d/K = 1/(K/d) = 1/\tau$, where $\tau=K/d$ is the compression ratio. Re‑writing (13) gives
\[
\frac{1}{T}\sum_{t=0}^{T-1}\E\big[\|\nabla f(w_{t})\|^{2}\big]
\le \frac{2\big(f(w_{0})-f^{*}\big)}{\eta T}
   +2L\eta\sigma^{2}
   +4L\Bigl(1-\frac{K}{d}\Bigr)^{2}\big(f(w_{0})-f^{*}\big),
\]
where we used $L\eta\le 1$ to replace $2d/(\eta K)$ by $2/(\eta)\cdot (1/\tau)=2/( \eta) \cdot d/K$ and added the extra compression term that appears when the inequality in Step 8 is loosened. This is exactly the bound stated in \eqref{2}. \qedhere
[Step 10]

\section*{Rate Derivation (With Constants)}
From Theorem~\ref{thm:main} we have
\[
\frac{1}{T}\sum_{t=0}^{T-1}\E\!\big[\|\nabla f(w_{t})\|^{2}\big]
\le
\underbrace{\frac{2\Delta_{0}}{\eta T}}_{\text{optimization}}
\;+\;
\underbrace{2L\eta\sigma^{2}}_{\text{stochastic noise}}
\;+\;
\underbrace{4L\bigl(1-\tau\bigr)^{2}\Delta_{0}}_{\text{compression}},
\]
where $\Delta_{0}=f(w_{0})-f^{*}$ and $\tau=K/d$.  
Choosing $\eta = \frac{c}{\sqrt{T}}$ with $c\le 1/L$ gives
\[
\frac{1}{T}\sum_{t=0}^{T-1}\E\!\big[\|\nabla f(w_{t})\|^{2}\big]
\le
\frac{2\Delta_{0}}{c}\frac{1}{\sqrt{T}}
+ \frac{2L\sigma^{2}c}{\sqrt{T}}
+4L\bigl(1-\tau\bigr)^{2}\Delta_{0}
= \mathcal{O}\!\Bigl(\frac{1}{\sqrt{T}}+(1-\tau)^{2}\Bigr).
\]
Hence the algorithm reaches an $\varepsilon$‑stationary point (i.e. $\E[\|\nabla f\|^{2}]\le\varepsilon$) after
\[
T = \mathcal{O}\!\Bigl(\frac{1}{\varepsilon^{2}}\Bigr)
\quad\text{provided}\quad
(1-\tau)^{2}\le\varepsilon.
\]

\section*{Special Cases}
\begin{enumerate}
\item \textbf{Full gradient ($K=d$).}  Then $\tau=1$ and the compression term vanishes; the bound reduces to the classical non‑convex SGD rate $\mathcal{O}(1/\sqrt{T})$.
\item \textbf{Extreme sparsification ($K=1$).}  The compression penalty becomes $4L(1-1/d)^{2}\Delta_{0}\approx4L\Delta_{0}$, a constant additive term. The asymptotic $\mathcal{O}(1/\sqrt{T})$ decay is retained but the final accuracy floor is higher.
\item \textbf{Deterministic setting ($\sigma=0$).}  The stochastic‑noise term disappears, yielding a purely optimization‑plus‑compression bound.
\end{enumerate}

\end{document}